\par\Needspace{6\baselineskip}\medskip
\begin{center}\hypertarget{native-174-article127}{}\pdfbookmark[2]{127.}{article-127}127.\end{center}
Once the three heliocentric positions have been found by the method described in the preceding
article, we can also proceed as follows. By solving the problem treated in Articles 85–105,
determine the remaining elements first from the first and second positions with their corresponding
time interval, and then in the same way from the second and third positions with their corresponding
time interval. This gives two values for each element, and any two of the differences between these
values may be taken for $X$ and $Y$. A considerable advantage of this method is that in the initial
hypotheses we may omit all the remaining elements except the two chosen to determine $X$ and $Y$.
The others need only be determined in the final calculation, based on the corrected values of $x,y$,
using just the first combination, just the second, or, as is generally preferable, the combination
of the first and third positions. The choice of the two elements is generally arbitrary and allows a
great variety of solutions. For example, we may pair the logarithm of the semiparameter with the
logarithm of the semimajor axis, either of those logarithms with the eccentricity, or any one of
these three elements with the longitude of perihelion. If the orbit turns out to be elliptic, we may
also pair any one of these four elements with the eccentric anomaly corresponding to the middle
position in each calculation. Formulas 27–30 in Article 96 then make the calculation particularly
straightforward. In special cases this choice calls for some care. For example, in orbits that
approach a parabola, the semimajor axis $a$ or its logarithm would be less suitable, since their
excessive variations could hardly be regarded as proportional to the variations in $x,y$. In such a
case it would be more appropriate to choose $\dfrac{1}{a}\mathpunct{\cdot}$ We need spend little time on these precautions,
since the fifth method, to be explained in the next article, surpasses the four presented so far in
almost every case.
